% Context from chapters-tex/sec-stochastic-optim.tex; for mathematical reference, not compilation.
\section{Stochastic Optimization}
\label{sec-stochastic-optim}

We study stochastic gradient methods for large datasets of size $n$ and for objectives defined as expectations.

                                                                                                                                                                  

                                                        
\subsection{Minimizing Sums and Expectations}

Many learning objectives are finite averages:
\eql{\label{eq-min-sums}
	\umin{x \in \RR^p} f(x) \eqdef \frac{1}{n} \sum_{i=1}^n f_i(x)
}
or expectations:
\eql{\label{eq-min-int}
	\umin{x \in \RR^p}  f(x) \eqdef \EE_{\zp \sim \pi}( f(x,\zp) ) = \int_{\Zz} f(x,z) \d\pi(z).
}
A finite average is an expectation under the uniform measure $\pi=\frac1n\sum_{i=1}^n\de_i$, with $f(x,i)=f_i(x)$. Conversely, drawing $(z_i)_i$ i.i.d. from $\pi$ and setting $f_i(x)=f(x,z_i)$ gives an empirical approximation of~\eqref{eq-min-int}. Under integrability, the empirical objective converges pointwise by the law of large numbers. Convergence of minima or minimizers requires additional assumptions.

For example, empirical risk minimization with linear predictors uses
\eql{\label{eq-stochastic-erm}
	f_i(x) = \ell(\dotp{a_i}{x},y_i)
	\qandq
	f(x,z) = \ell(\dotp{a}{x},y)
}
for $z=(a,y) \in \Zz = (\Aa=\RR^p) \times \Yy$ (typically $\Yy=\RR$ or $\Yy=\{-1,+1\}$ for regression and classification), where $\ell$ is some loss function. 
 
We illustrate the methods with binary logistic classification, using
\eql{\label{eq-stoch-logistic}
	\loss(s,y) \eqdef \log( 1+\exp(-sy) ).
}
The same computational approach extends to general parametric models, including deep neural networks.

The stochastic methods below apply both to finite sums~\eqref{eq-min-sums} and to expectations~\eqref{eq-min-int}, under the corresponding unbiasedness and moment assumptions. We present them for finite sums. When $n$ is large, the aim is to make each iteration inexpensive without evaluating every summand.

When the summands $f_i(x)$ are similar, a single gradient can approximate the average well: $\nabla f_i \approx \nabla f$. Evaluating $\nabla f_i$ is typically $n$ times cheaper than forming the full gradient.
 
Stochastic optimization is nevertheless not always faster than batch gradient descent.
 
For moderate $n$, a deterministic method may use the available computation more effectively.
 
When $n$ is too large for a full pass to be practical, stochastic methods divide the work into small updates and begin improving the model immediately. Mini-batches also permit parallel computation.


                                                        
\subsection{Batch Gradient Descent (BGD)}

Batch gradient descent (BGD), analyzed in Section~\ref{sec-grad-desc-basic}, iterates
\eq{
	x_{k+1} = x_k - \tau_k \nabla f(x_k)
}
For a convex objective with an $L$-Lipschitz gradient and a minimizer, the condition $0<\tau_{\min}\leq\tau_k\leq\tau_{\max}<2/L$ ensures convergence. Strong convexity strengthens this to a linear rate.

For a finite-sum objective, the gradient is
\eql{\label{eq-full-grad}
	\nabla f(x) = \frac{1}{n} \sum_{i=1}^n \nabla f_i(x)
}
Its cost is typically $O(np)$ when each $\nabla f_i$ can be evaluated in time linear in $p$.

For an ERM summand of the form~\eqref{eq-stochastic-erm}, Taylor expansion gives
\begin{align*}
	f_i(x+\epsilon) &= \ell(\dotp{a_i}{x} + \dotp{a_i}{\epsilon},y_i)
	= \ell(\dotp{a_i}{x},y_i) + \ell'(\dotp{a_i}{x},y_i)\dotp{a_i}{\epsilon} + o(\norm{\epsilon}) \\
	&= f_i(x) + \dotp{\ell'(\dotp{a_i}{x},y_i)a_i}{\epsilon} + o(\norm{\epsilon}),
\end{align*}
Here $\ell'(s,y)$ denotes differentiation with respect to the score $s$, with the response $y$ fixed. Consequently,
\eql{\label{eq-grad-formula}
	\nabla f_i(x) =  \ell'( \dotp{a_i}{x},y_i ) a_i.
}
For the logistic loss, one has 
\eq{
	\loss'(s,y)=-y \frac{e^{-sy}}{ 1+e^{-sy} }.
}



\begin{figure}
\centering
\begin{tabular}{cc}
\includegraphics[width=.45\linewidth]{ml/sgd/error-bgd-1} &
\includegraphics[width=.45\linewidth]{ml/sgd/error-bgd-2} \\
$f(x_k)$ & $\log_{10}(f(x_k)-f(x^\star))$ 
\end{tabular}
\caption{\label{fig-bgd}
Objective-error evolution for batch gradient descent in logistic classification.
}
\end{figure}



                                                        
\subsection{Stochastic Gradient Descent (SGD)}

\begin{figure}[tbp]
\centering
\BookPanel{.485}{1}{tikz/advanced-unbiased-gradient}\hfill
\BookPanel{.485}{2}{tikz/advanced-sgd-trajectory}
\caption{Unbiased stochastic gradient estimate (panel 1) and a schematic SGD trajectory (panel 2).}
\label{fig-review-advanced-unbiased-gradient}
\label{fig-review-advanced-sgd-trajectory}
\end{figure}
For very large $n$, computing the full gradient $\nabla f$ as in~\eqref{eq-full-grad} is prohibitive.  
 
SGD replaces the full gradient by the gradient of one summand $f_{i}$, with $i$ sampled uniformly. This gives an unbiased estimate:
\eql{\label{eq-unbiased-grad}
	\EE_{\ip}{ \nabla f_{\ip}(x) } = \nabla f(x)
}
where $\ip$ is a random variable distributed uniformly in $\{1,\ldots,n\}$.


Starting from $x_0$, stochastic gradient descent (SGD) iterates
\eq{
	x_{k+1} = x_k - \tau_k \nabla f_{i(k)}(x_k)
}
where, for each iteration index $k$, $i(k)$
is drawn independently of previous draws, uniformly from $\{1,\ldots,n\}$.
 
The iterates $x_{k+1}$ are random. We therefore study convergence to a minimizer of $f$ in a probabilistic sense, such as mean square or probability, together with the corresponding rates.

A batch-gradient step costs $O(np)$,
whereas an SGD step costs $O(p)$. This makes SGD attractive
when $n$ is large and repeated passes through the data are costly.
Redundancy between observations can sometimes yield
accurate predictions after only $k \ll n$ updates, before all
observations have been used.

The step-size schedule $\tau_k$ must balance two requirements: it
usually needs to decrease to suppress stochastic noise,
but not so quickly that the method stops making progress
before reaching a minimizer.


A common schedule has asymptotic behavior $\tau_k \sim k^{-1}$
as $k\rightarrow +\infty$. One such schedule is
\eql{\label{eq-stepsize-sgd}
	\tau_k \eqdef \frac{\tau_0}{1 + k/k_0}
}
where $\tau_0>0$ sets the in